\documentclass[10pt]{article}
\usepackage[utf8]{inputenc}
\usepackage[T1]{fontenc}
\usepackage{amsmath}
\usepackage{amsfonts}
\usepackage{amssymb}
\usepackage[version=4]{mhchem}
\usepackage{stmaryrd}
\usepackage{bbold}
\usepackage{caption}
\usepackage{graphicx}
\usepackage[export]{adjustbox}
\graphicspath{ {./images/} }

\begin{document}
\captionsetup{singlelinecheck=false}
Second Welfare Theorem, an application

\section*{Negishi's approach}
\begin{itemize}
  \item Use the welfare theorems to identify WE.
  \item Analysis is in utility space, not commodity space.
  \item Useful with many goods and few agents.
  \item Proposed by Negishi (Negishi 1960).
  \item Employed by Arrow and Hahn (Arrow and Hahn 1971) and many others.
  \item Overview in Macolell-etal, (Mas-Colell, Whinston, and Green 1995, "16.E, 17.Appendix A")
\end{itemize}

\section*{Preliminaries}
\begin{itemize}
  \item Let $\mathcal{E}$ be an Arrow-Debreu economy.
  \item $(\forall i)\left(\exists u_{i}, u_{i}: X_{i} \rightarrow \mathbb{R}\right)$ such that
\end{itemize}

$$
\left(\forall x_{i}, x_{i}^{\prime} \in X_{i}\right) \quad\left[x_{i} \succeq_{i} x_{i}^{\prime} \Longleftrightarrow u_{i}\left(x_{i}\right) \geq u_{i}\left(x_{i}^{\prime}\right)\right] .
$$

\begin{itemize}
  \item Let $x=\left(x_{1}, \ldots, x_{i}, \ldots, x_{I}\right)$.
  \item $u(x)$ denotes the vector $\left(u_{1}\left(x_{1}\right), \ldots, u_{I}\left(x_{I}\right)\right)$.
  \item Let $u: \mathbb{R}_{+}^{L I} \rightarrow \mathbb{R}^{I}$ be the map $x \mapsto\left(u_{1}\left(x_{1}\right), \ldots, u_{I}\left(x_{I}\right)\right)$.
\end{itemize}

\section*{Utility possibility set}
Definition 1 The utility possibility set is

$$
U=\left\{u \in \mathbb{R}^{I}: \exists \text { an allocation } x \text { with } u \leq u(x)\right\}
$$

and its Pareto frontier is

$$
P=\left\{u \in U: \nexists u^{\prime} \in U, u^{\prime}>u\right\} .
$$

An allocation $x^{*}$ is Pareto optimal if and only if $u\left(x^{*}\right) \in P$.\\
Similar definitions apply to economies with production.

\section*{Utility possibility set}
\begin{figure}[h]
\begin{center}
\captionsetup{labelformat=empty}
\caption{Utility set: utilities normalized}
  \includegraphics[alt={},max width=\textwidth]{502add65-2db5-4878-b75f-d85c1cc8d7d7-06_1430_2080_425_166}
\end{center}
\end{figure}

\section*{Social welfare function}
\begin{itemize}
  \item Suppose society's values are summarized by a function
\end{itemize}

$$
W(u)=W\left(u_{1}, \ldots, u_{I}\right) .
$$

\begin{itemize}
  \item Pick $a \in \mathbb{R}^{I}$ with $(\forall i) a_{i} \geq 0$. Define
\end{itemize}

$$
W(u)=\sum_{i=1}^{I} a_{i} u_{i}=a \cdot u
$$

\section*{Lemmas}
Lemma 1 If $\exists a \in \mathbb{R}_{++}^{I}$ and $u^{*} \in U$ such that $(\forall u \in U) a \cdot u^{*} \geq a \cdot u$, then $u^{*} \in P$.

Lemma 2 If $U$ is convex, then for any $u^{*} \in P,\left(\exists a \in \mathbb{R}_{+}^{I}\right) a \neq 0$ and $(\forall u \in U) a \cdot u^{*} \geq a \cdot u$.

\section*{Lemma 1 and Lemma 2}
\includegraphics[max width=\textwidth, alt={}, center]{502add65-2db5-4878-b75f-d85c1cc8d7d7-09_1352_1269_389_157}\\
\includegraphics[max width=\textwidth, alt={}, center]{502add65-2db5-4878-b75f-d85c1cc8d7d7-09_1350_1267_387_1456}

\section*{Stylized method}
\begin{itemize}
  \item Let $\Delta=\left\{s \in \mathbb{R}_{+}^{I}: \sum_{i} s_{i}=1\right\}$.
  \item Fix $a \in \Delta$, pick $u(a) \in U$ so that $u(a) \in P$.
  \item $u(a)$ is derived from an allocation $(x(a), y(a))$.
  \item The Weak Second Welfare Theorem implies
  \item $\exists p(a), y_{j}(a)$ maximizes firm $j$ 's profit and,
  \item $x_{i}(a)$ is in $i$ 's quasi demand.
  \item Identify transfer: $(\forall i), t_{i}=p(a) \cdot\left[\omega_{i}-x_{i}(a)\right]$.
  \item If $(\exists a)$ such that $(\forall i) t_{i}=0$. This identifies a WE.
\end{itemize}

Note Arrow's exceptional economy.

\section*{References}
Arrow, Kenneth, and Frank H. Hahn. 1971. General Competitive Analysis. 1st ed. Vol. 12. North Holland.\\
Mas-Colell, Andreu, Michael Whinston, and Jerry Green. 1995. Microeconomic Theory. New York, Oxford: Oxford University Press.\\
Negishi, T. 1960. "Welfare Economics and the Existence of Equilibrium for a Competitive Economy."\\
Metroeconomica 12: 92-97.


\end{document}